\chapter{Refactoring Plans}

% === Source file: refactor/clawnet.md ===
\section{Clawnet refactor (protocol + auth unification)}


\subsection{Hi}


Hi Peter — great direction; this unlocks simpler UX + stronger security.

\subsection{Purpose}


Single, rigorous document for:

\begin{itemize}
  \item Current state: protocols, flows, trust boundaries.
  \item Pain points: approvals, multi‑hop routing, UI duplication.
  \item Proposed new state: one protocol, scoped roles, unified auth/pairing, TLS pinning.
  \item Identity model: stable IDs + cute slugs.
  \item Migration plan, risks, open questions.
\end{itemize}


\subsection{Goals (from discussion)}


\begin{itemize}
  \item One protocol for all clients (mac app, CLI, iOS, Android, headless node).
  \item Every network participant authenticated + paired.
  \item Role clarity: nodes vs operators.
  \item Central approvals routed to where the user is.
  \item TLS encryption + optional pinning for all remote traffic.
  \item Minimal code duplication.
  \item Single machine should appear once (no UI/node duplicate entry).
\end{itemize}


\subsection{Non‑goals (explicit)}


\begin{itemize}
  \item Remove capability separation (still need least‑privilege).
  \item Expose full gateway control plane without scope checks.
  \item Make auth depend on human labels (slugs remain non‑security).
\end{itemize}


\hrulefill


\section{Current state (as‑is)}


\subsection{Two protocols}


\subsubsection{1) Gateway WebSocket (control plane)}


\begin{itemize}
  \item Full API surface: config, channels, models, sessions, agent runs, logs, nodes, etc.
  \item Default bind: loopback. Remote access via SSH/Tailscale.
  \item Auth: token/password via \texttt{connect}.
  \item No TLS pinning (relies on loopback/tunnel).
  \item Code:
  \item \texttt{src/gateway/server/ws-connection/message-handler.ts}
  \item \texttt{src/gateway/client.ts}
  \item \texttt{docs/gateway/protocol.md}
\end{itemize}


\subsubsection{2) Bridge (node transport)}


\begin{itemize}
  \item Narrow allowlist surface, node identity + pairing.
  \item JSONL over TCP; optional TLS + cert fingerprint pinning.
  \item TLS advertises fingerprint in discovery TXT.
  \item Code:
  \item \texttt{src/infra/bridge/server/connection.ts}
  \item \texttt{src/gateway/server-bridge.ts}
  \item \texttt{src/node-host/bridge-client.ts}
  \item \texttt{docs/gateway/bridge-protocol.md}
\end{itemize}


\subsection{Control plane clients today}


\begin{itemize}
  \item CLI → Gateway WS via \texttt{callGateway} (\texttt{src/gateway/call.ts}).
  \item macOS app UI → Gateway WS (\texttt{GatewayConnection}).
  \item Web Control UI → Gateway WS.
  \item ACP → Gateway WS.
  \item Browser control uses its own HTTP control server.
\end{itemize}


\subsection{Nodes today}


\begin{itemize}
  \item macOS app in node mode connects to Gateway bridge (\texttt{MacNodeBridgeSession}).
  \item iOS/Android apps connect to Gateway bridge.
  \item Pairing + per‑node token stored on gateway.
\end{itemize}


\subsection{Current approval flow (exec)}


\begin{itemize}
  \item Agent uses \texttt{system.run} via Gateway.
  \item Gateway invokes node over bridge.
  \item Node runtime decides approval.
  \item UI prompt shown by mac app (when node == mac app).
  \item Node returns \texttt{invoke-res} to Gateway.
  \item Multi‑hop, UI tied to node host.
\end{itemize}


\subsection{Presence + identity today}


\begin{itemize}
  \item Gateway presence entries from WS clients.
  \item Node presence entries from bridge.
  \item mac app can show two entries for same machine (UI + node).
  \item Node identity stored in pairing store; UI identity separate.
\end{itemize}


\hrulefill


\section{Problems / pain points}


\begin{itemize}
  \item Two protocol stacks to maintain (WS + Bridge).
  \item Approvals on remote nodes: prompt appears on node host, not where user is.
  \item TLS pinning only exists for bridge; WS depends on SSH/Tailscale.
  \item Identity duplication: same machine shows as multiple instances.
  \item Ambiguous roles: UI + node + CLI capabilities not clearly separated.
\end{itemize}


\hrulefill


\section{Proposed new state (Clawnet)}


\subsection{One protocol, two roles}


Single WS protocol with role + scope.

\begin{itemize}
  \item \textbf{Role: node} (capability host)
  \item \textbf{Role: operator} (control plane)
  \item Optional \textbf{scope} for operator:
  \item \texttt{operator.read} (status + viewing)
  \item \texttt{operator.write} (agent run, sends)
  \item \texttt{operator.admin} (config, channels, models)
\end{itemize}


\subsubsection{Role behaviors}


\textbf{Node}

\begin{itemize}
  \item Can register capabilities (\texttt{caps}, \texttt{commands}, permissions).
  \item Can receive \texttt{invoke} commands (\texttt{system.run}, \texttt{camera.\textit{}, \texttt{canvas.}}, \texttt{screen.record}, etc).
  \item Can send events: \texttt{voice.transcript}, \texttt{agent.request}, \texttt{chat.subscribe}.
  \item Cannot call config/models/channels/sessions/agent control plane APIs.
\end{itemize}


\textbf{Operator}

\begin{itemize}
  \item Full control plane API, gated by scope.
  \item Receives all approvals.
  \item Does not directly execute OS actions; routes to nodes.
\end{itemize}


\subsubsection{Key rule}


Role is per‑connection, not per device. A device may open both roles, separately.

\hrulefill


\section{Unified authentication + pairing}


\subsection{Client identity}


Every client provides:

\begin{itemize}
  \item \texttt{deviceId} (stable, derived from device key).
  \item \texttt{displayName} (human name).
  \item \texttt{role} + \texttt{scope} + \texttt{caps} + \texttt{commands}.
\end{itemize}


\subsection{Pairing flow (unified)}


\begin{itemize}
  \item Client connects unauthenticated.
  \item Gateway creates a \textbf{pairing request} for that \texttt{deviceId}.
  \item Operator receives prompt; approves/denies.
  \item Gateway issues credentials bound to:
  \item device public key
  \item role(s)
  \item scope(s)
  \item capabilities/commands
  \item Client persists token, reconnects authenticated.
\end{itemize}


\subsection{Device‑bound auth (avoid bearer token replay)}


Preferred: device keypairs.

\begin{itemize}
  \item Device generates keypair once.
  \item \texttt{deviceId = fingerprint(publicKey)}.
  \item Gateway sends nonce; device signs; gateway verifies.
  \item Tokens are issued to a public key (proof‑of‑possession), not a string.
\end{itemize}


Alternatives:

\begin{itemize}
  \item mTLS (client certs): strongest, more ops complexity.
  \item Short‑lived bearer tokens only as a temporary phase (rotate + revoke early).
\end{itemize}


\subsection{Silent approval (SSH heuristic)}


Define it precisely to avoid a weak link. Prefer one:

\begin{itemize}
  \item \textbf{Local‑only}: auto‑pair when client connects via loopback/Unix socket.
  \item \textbf{Challenge via SSH}: gateway issues nonce; client proves SSH by fetching it.
  \item \textbf{Physical presence window}: after a local approval on gateway host UI, allow auto‑pair for a short window (e.g. 10 minutes).
\end{itemize}


Always log + record auto‑approvals.

\hrulefill


\section{TLS everywhere (dev + prod)}


\subsection{Reuse existing bridge TLS}


Use current TLS runtime + fingerprint pinning:

\begin{itemize}
  \item \texttt{src/infra/bridge/server/tls.ts}
  \item fingerprint verification logic in \texttt{src/node-host/bridge-client.ts}
\end{itemize}


\subsection{Apply to WS}


\begin{itemize}
  \item WS server supports TLS with same cert/key + fingerprint.
  \item WS clients can pin fingerprint (optional).
  \item Discovery advertises TLS + fingerprint for all endpoints.
  \item Discovery is locator hints only; never a trust anchor.
\end{itemize}


\subsection{Why}


\begin{itemize}
  \item Reduce reliance on SSH/Tailscale for confidentiality.
  \item Make remote mobile connections safe by default.
\end{itemize}


\hrulefill


\section{Approvals redesign (centralized)}


\subsection{Current}


Approval happens on node host (mac app node runtime). Prompt appears where node runs.

\subsection{Proposed}


Approval is \textbf{gateway‑hosted}, UI delivered to operator clients.

\subsubsection{New flow}


\begin{enumerate}
  \item Gateway receives \texttt{system.run} intent (agent).
  \item Gateway creates approval record: \texttt{approval.requested}.
  \item Operator UI(s) show prompt.
  \item Approval decision sent to gateway: \texttt{approval.resolve}.
  \item Gateway invokes node command if approved.
  \item Node executes, returns \texttt{invoke-res}.
\end{enumerate}


\subsubsection{Approval semantics (hardening)}


\begin{itemize}
  \item Broadcast to all operators; only the active UI shows a modal (others get a toast).
  \item First resolution wins; gateway rejects subsequent resolves as already settled.
  \item Default timeout: deny after N seconds (e.g. 60s), log reason.
  \item Resolution requires \texttt{operator.approvals} scope.
\end{itemize}


\subsection{Benefits}


\begin{itemize}
  \item Prompt appears where user is (mac/phone).
  \item Consistent approvals for remote nodes.
  \item Node runtime stays headless; no UI dependency.
\end{itemize}


\hrulefill


\section{Role clarity examples}


\subsection{iPhone app}


\begin{itemize}
  \item \textbf{Node role} for: mic, camera, voice chat, location, push‑to‑talk.
  \item Optional \textbf{operator.read} for status and chat view.
  \item Optional \textbf{operator.write/admin} only when explicitly enabled.
\end{itemize}


\subsection{macOS app}


\begin{itemize}
  \item Operator role by default (control UI).
  \item Node role when “Mac node” enabled (system.run, screen, camera).
  \item Same deviceId for both connections → merged UI entry.
\end{itemize}


\subsection{CLI}


\begin{itemize}
  \item Operator role always.
  \item Scope derived by subcommand:
  \item \texttt{status}, \texttt{logs} → read
  \item \texttt{agent}, \texttt{message} → write
  \item \texttt{config}, \texttt{channels} → admin
  \item approvals + pairing → \texttt{operator.approvals} / \texttt{operator.pairing}
\end{itemize}


\hrulefill


\section{Identity + slugs}


\subsection{Stable ID}


Required for auth; never changes.
Preferred:

\begin{itemize}
  \item Keypair fingerprint (public key hash).
\end{itemize}


\subsection{Cute slug (lobster‑themed)}


Human label only.

\begin{itemize}
  \item Example: \texttt{scarlet-claw}, \texttt{saltwave}, \texttt{mantis-pinch}.
  \item Stored in gateway registry, editable.
  \item Collision handling: \texttt{-2}, \texttt{-3}.
\end{itemize}


\subsection{UI grouping}


Same \texttt{deviceId} across roles → single “Instance” row:

\begin{itemize}
  \item Badge: \texttt{operator}, \texttt{node}.
  \item Shows capabilities + last seen.
\end{itemize}


\hrulefill


\section{Migration strategy}


\subsection{Phase 0: Document + align}


\begin{itemize}
  \item Publish this doc.
  \item Inventory all protocol calls + approval flows.
\end{itemize}


\subsection{Phase 1: Add roles/scopes to WS}


\begin{itemize}
  \item Extend \texttt{connect} params with \texttt{role}, \texttt{scope}, \texttt{deviceId}.
  \item Add allowlist gating for node role.
\end{itemize}


\subsection{Phase 2: Bridge compatibility}


\begin{itemize}
  \item Keep bridge running.
  \item Add WS node support in parallel.
  \item Gate features behind config flag.
\end{itemize}


\subsection{Phase 3: Central approvals}


\begin{itemize}
  \item Add approval request + resolve events in WS.
  \item Update mac app UI to prompt + respond.
  \item Node runtime stops prompting UI.
\end{itemize}


\subsection{Phase 4: TLS unification}


\begin{itemize}
  \item Add TLS config for WS using bridge TLS runtime.
  \item Add pinning to clients.
\end{itemize}


\subsection{Phase 5: Deprecate bridge}


\begin{itemize}
  \item Migrate iOS/Android/mac node to WS.
  \item Keep bridge as fallback; remove once stable.
\end{itemize}


\subsection{Phase 6: Device‑bound auth}


\begin{itemize}
  \item Require key‑based identity for all non‑local connections.
  \item Add revocation + rotation UI.
\end{itemize}


\hrulefill


\section{Security notes}


\begin{itemize}
  \item Role/allowlist enforced at gateway boundary.
  \item No client gets “full” API without operator scope.
  \item Pairing required for \_all\_ connections.
  \item TLS + pinning reduces MITM risk for mobile.
  \item SSH silent approval is a convenience; still recorded + revocable.
  \item Discovery is never a trust anchor.
  \item Capability claims are verified against server allowlists by platform/type.
\end{itemize}


\section{Streaming + large payloads (node media)}


WS control plane is fine for small messages, but nodes also do:

\begin{itemize}
  \item camera clips
  \item screen recordings
  \item audio streams
\end{itemize}


Options:

\begin{enumerate}
  \item WS binary frames + chunking + backpressure rules.
  \item Separate streaming endpoint (still TLS + auth).
  \item Keep bridge longer for media‑heavy commands, migrate last.
\end{enumerate}


Pick one before implementation to avoid drift.

\section{Capability + command policy}


\begin{itemize}
  \item Node‑reported caps/commands are treated as \textbf{claims}.
  \item Gateway enforces per‑platform allowlists.
  \item Any new command requires operator approval or explicit allowlist change.
  \item Audit changes with timestamps.
\end{itemize}


\section{Audit + rate limiting}


\begin{itemize}
  \item Log: pairing requests, approvals/denials, token issuance/rotation/revocation.
  \item Rate‑limit pairing spam and approval prompts.
\end{itemize}


\section{Protocol hygiene}


\begin{itemize}
  \item Explicit protocol version + error codes.
  \item Reconnect rules + heartbeat policy.
  \item Presence TTL and last‑seen semantics.
\end{itemize}


\hrulefill


\section{Open questions}


\begin{enumerate}
  \item Single device running both roles: token model
\end{enumerate}

\begin{itemize}
  \item Recommend separate tokens per role (node vs operator).
  \item Same deviceId; different scopes; clearer revocation.
\end{itemize}


\begin{enumerate}
  \item Operator scope granularity
\end{enumerate}

\begin{itemize}
  \item read/write/admin + approvals + pairing (minimum viable).
  \item Consider per‑feature scopes later.
\end{itemize}


\begin{enumerate}
  \item Token rotation + revocation UX
\end{enumerate}

\begin{itemize}
  \item Auto‑rotate on role change.
  \item UI to revoke by deviceId + role.
\end{itemize}


\begin{enumerate}
  \item Discovery
\end{enumerate}

\begin{itemize}
  \item Extend current Bonjour TXT to include WS TLS fingerprint + role hints.
  \item Treat as locator hints only.
\end{itemize}


\begin{enumerate}
  \item Cross‑network approval
\end{enumerate}

\begin{itemize}
  \item Broadcast to all operator clients; active UI shows modal.
  \item First response wins; gateway enforces atomicity.
\end{itemize}


\hrulefill


\section{Summary (TL;DR)}


\begin{itemize}
  \item Today: WS control plane + Bridge node transport.
  \item Pain: approvals + duplication + two stacks.
  \item Proposal: one WS protocol with explicit roles + scopes, unified pairing + TLS pinning, gateway‑hosted approvals, stable device IDs + cute slugs.
  \item Outcome: simpler UX, stronger security, less duplication, better mobile routing.
\end{itemize}



\clearpage

% === Source file: refactor/exec-host.md ===
\section{Exec host refactor plan}


\subsection{Goals}


\begin{itemize}
  \item Add \texttt{exec.host} + \texttt{exec.security} to route execution across \textbf{sandbox}, \textbf{gateway}, and \textbf{node}.
  \item Keep defaults \textbf{safe}: no cross-host execution unless explicitly enabled.
  \item Split execution into a \textbf{headless runner service} with optional UI (macOS app) via local IPC.
  \item Provide \textbf{per-agent} policy, allowlist, ask mode, and node binding.
  \item Support \textbf{ask modes} that work \_with\_ or \_without\_ allowlists.
  \item Cross-platform: Unix socket + token auth (macOS/Linux/Windows parity).
\end{itemize}


\subsection{Non-goals}


\begin{itemize}
  \item No legacy allowlist migration or legacy schema support.
  \item No PTY/streaming for node exec (aggregated output only).
  \item No new network layer beyond the existing Bridge + Gateway.
\end{itemize}


\subsection{Decisions (locked)}


\begin{itemize}
  \item \textbf{Config keys:} \texttt{exec.host} + \texttt{exec.security} (per-agent override allowed).
  \item \textbf{Elevation:} keep \texttt{/elevated} as an alias for gateway full access.
  \item \textbf{Ask default:} \texttt{on-miss}.
  \item \textbf{Approvals store:} \texttt{\textasciitilde{}/.openclaw/exec-approvals.json} (JSON, no legacy migration).
  \item \textbf{Runner:} headless system service; UI app hosts a Unix socket for approvals.
  \item \textbf{Node identity:} use existing \texttt{nodeId}.
  \item \textbf{Socket auth:} Unix socket + token (cross-platform); split later if needed.
  \item \textbf{Node host state:} \texttt{\textasciitilde{}/.openclaw/node.json} (node id + pairing token).
  \item \textbf{macOS exec host:} run \texttt{system.run} inside the macOS app; node host service forwards requests over local IPC.
  \item \textbf{No XPC helper:} stick to Unix socket + token + peer checks.
\end{itemize}


\subsection{Key concepts}


\subsubsection{Host}


\begin{itemize}
  \item \texttt{sandbox}: Docker exec (current behavior).
  \item \texttt{gateway}: exec on gateway host.
  \item \texttt{node}: exec on node runner via Bridge (\texttt{system.run}).
\end{itemize}


\subsubsection{Security mode}


\begin{itemize}
  \item \texttt{deny}: always block.
  \item \texttt{allowlist}: allow only matches.
  \item \texttt{full}: allow everything (equivalent to elevated).
\end{itemize}


\subsubsection{Ask mode}


\begin{itemize}
  \item \texttt{off}: never ask.
  \item \texttt{on-miss}: ask only when allowlist does not match.
  \item \texttt{always}: ask every time.
\end{itemize}


Ask is \textbf{independent} of allowlist; allowlist can be used with \texttt{always} or \texttt{on-miss}.

\subsubsection{Policy resolution (per exec)}


\begin{enumerate}
  \item Resolve \texttt{exec.host} (tool param → agent override → global default).
  \item Resolve \texttt{exec.security} and \texttt{exec.ask} (same precedence).
  \item If host is \texttt{sandbox}, proceed with local sandbox exec.
  \item If host is \texttt{gateway} or \texttt{node}, apply security + ask policy on that host.
\end{enumerate}


\subsection{Default safety}


\begin{itemize}
  \item Default \texttt{exec.host = sandbox}.
  \item Default \texttt{exec.security = deny} for \texttt{gateway} and \texttt{node}.
  \item Default \texttt{exec.ask = on-miss} (only relevant if security allows).
  \item If no node binding is set, \textbf{agent may target any node}, but only if policy allows it.
\end{itemize}


\subsection{Config surface}


\subsubsection{Tool parameters}


\begin{itemize}
  \item \texttt{exec.host} (optional): \texttt{sandbox | gateway | node}.
  \item \texttt{exec.security} (optional): \texttt{deny | allowlist | full}.
  \item \texttt{exec.ask} (optional): \texttt{off | on-miss | always}.
  \item \texttt{exec.node} (optional): node id/name to use when \texttt{host=node}.
\end{itemize}


\subsubsection{Config keys (global)}


\begin{itemize}
  \item \texttt{tools.exec.host}
  \item \texttt{tools.exec.security}
  \item \texttt{tools.exec.ask}
  \item \texttt{tools.exec.node} (default node binding)
\end{itemize}


\subsubsection{Config keys (per agent)}


\begin{itemize}
  \item \texttt{agents.list[].tools.exec.host}
  \item \texttt{agents.list[].tools.exec.security}
  \item \texttt{agents.list[].tools.exec.ask}
  \item \texttt{agents.list[].tools.exec.node}
\end{itemize}


\subsubsection{Alias}


\begin{itemize}
  \item \texttt{/elevated on} = set \texttt{tools.exec.host=gateway}, \texttt{tools.exec.security=full} for the agent session.
  \item \texttt{/elevated off} = restore previous exec settings for the agent session.
\end{itemize}


\subsection{Approvals store (JSON)}


Path: \texttt{\textasciitilde{}/.openclaw/exec-approvals.json}

Purpose:

\begin{itemize}
  \item Local policy + allowlists for the \textbf{execution host} (gateway or node runner).
  \item Ask fallback when no UI is available.
  \item IPC credentials for UI clients.
\end{itemize}


Proposed schema (v1):

\begin{lstlisting}[language=json]
{
  "version": 1,
  "socket": {
    "path": "~/.openclaw/exec-approvals.sock",
    "token": "base64-opaque-token"
  },
  "defaults": {
    "security": "deny",
    "ask": "on-miss",
    "askFallback": "deny"
  },
  "agents": {
    "agent-id-1": {
      "security": "allowlist",
      "ask": "on-miss",
      "allowlist": [
        {
          "pattern": "~/Projects/**/bin/rg",
          "lastUsedAt": 0,
          "lastUsedCommand": "rg -n TODO",
          "lastResolvedPath": "/Users/user/Projects/.../bin/rg"
        }
      ]
    }
  }
}
\end{lstlisting}


Notes:

\begin{itemize}
  \item No legacy allowlist formats.
  \item \texttt{askFallback} applies only when \texttt{ask} is required and no UI is reachable.
  \item File permissions: \texttt{0600}.
\end{itemize}


\subsection{Runner service (headless)}


\subsubsection{Role}


\begin{itemize}
  \item Enforce \texttt{exec.security} + \texttt{exec.ask} locally.
  \item Execute system commands and return output.
  \item Emit Bridge events for exec lifecycle (optional but recommended).
\end{itemize}


\subsubsection{Service lifecycle}


\begin{itemize}
  \item Launchd/daemon on macOS; system service on Linux/Windows.
  \item Approvals JSON is local to the execution host.
  \item UI hosts a local Unix socket; runners connect on demand.
\end{itemize}


\subsection{UI integration (macOS app)}


\subsubsection{IPC}


\begin{itemize}
  \item Unix socket at \texttt{\textasciitilde{}/.openclaw/exec-approvals.sock} (0600).
  \item Token stored in \texttt{exec-approvals.json} (0600).
  \item Peer checks: same-UID only.
  \item Challenge/response: nonce + HMAC(token, request-hash) to prevent replay.
  \item Short TTL (e.g., 10s) + max payload + rate limit.
\end{itemize}


\subsubsection{Ask flow (macOS app exec host)}


\begin{enumerate}
  \item Node service receives \texttt{system.run} from gateway.
  \item Node service connects to the local socket and sends the prompt/exec request.
  \item App validates peer + token + HMAC + TTL, then shows dialog if needed.
  \item App executes the command in UI context and returns output.
  \item Node service returns output to gateway.
\end{enumerate}


If UI missing:

\begin{itemize}
  \item Apply \texttt{askFallback} (\texttt{deny|allowlist|full}).
\end{itemize}


\subsubsection{Diagram (SCI)}


\begin{lstlisting}
Agent -> Gateway -> Bridge -> Node Service (TS)
                         |  IPC (UDS + token + HMAC + TTL)
                         v
                     Mac App (UI + TCC + system.run)
\end{lstlisting}


\subsection{Node identity + binding}


\begin{itemize}
  \item Use existing \texttt{nodeId} from Bridge pairing.
  \item Binding model:
  \item \texttt{tools.exec.node} restricts the agent to a specific node.
  \item If unset, agent can pick any node (policy still enforces defaults).
  \item Node selection resolution:
  \item \texttt{nodeId} exact match
  \item \texttt{displayName} (normalized)
  \item \texttt{remoteIp}
  \item \texttt{nodeId} prefix (>= 6 chars)
\end{itemize}


\subsection{Eventing}


\subsubsection{Who sees events}


\begin{itemize}
  \item System events are \textbf{per session} and shown to the agent on the next prompt.
  \item Stored in the gateway in-memory queue (\texttt{enqueueSystemEvent}).
\end{itemize}


\subsubsection{Event text}


\begin{itemize}
  \item \texttt{Exec started (node=, id=)}
  \item \texttt{Exec finished (node=, id=, code=)} + optional output tail
  \item \texttt{Exec denied (node=, id=, )}
\end{itemize}


\subsubsection{Transport}


Option A (recommended):

\begin{itemize}
  \item Runner sends Bridge \texttt{event} frames \texttt{exec.started} / \texttt{exec.finished}.
  \item Gateway \texttt{handleBridgeEvent} maps these into \texttt{enqueueSystemEvent}.
\end{itemize}


Option B:

\begin{itemize}
  \item Gateway \texttt{exec} tool handles lifecycle directly (synchronous only).
\end{itemize}


\subsection{Exec flows}


\subsubsection{Sandbox host}


\begin{itemize}
  \item Existing \texttt{exec} behavior (Docker or host when unsandboxed).
  \item PTY supported in non-sandbox mode only.
\end{itemize}


\subsubsection{Gateway host}


\begin{itemize}
  \item Gateway process executes on its own machine.
  \item Enforces local \texttt{exec-approvals.json} (security/ask/allowlist).
\end{itemize}


\subsubsection{Node host}


\begin{itemize}
  \item Gateway calls \texttt{node.invoke} with \texttt{system.run}.
  \item Runner enforces local approvals.
  \item Runner returns aggregated stdout/stderr.
  \item Optional Bridge events for start/finish/deny.
\end{itemize}


\subsection{Output caps}


\begin{itemize}
  \item Cap combined stdout+stderr at \textbf{200k}; keep \textbf{tail 20k} for events.
  \item Truncate with a clear suffix (e.g., \texttt{"… (truncated)"}).
\end{itemize}


\subsection{Slash commands}


\begin{itemize}
  \item \texttt{/exec host= security= ask= node=}
  \item Per-agent, per-session overrides; non-persistent unless saved via config.
  \item \texttt{/elevated on|off|ask|full} remains a shortcut for \texttt{host=gateway security=full} (with \texttt{full} skipping approvals).
\end{itemize}


\subsection{Cross-platform story}


\begin{itemize}
  \item The runner service is the portable execution target.
  \item UI is optional; if missing, \texttt{askFallback} applies.
  \item Windows/Linux support the same approvals JSON + socket protocol.
\end{itemize}


\subsection{Implementation phases}


\subsubsection{Phase 1: config + exec routing}


\begin{itemize}
  \item Add config schema for \texttt{exec.host}, \texttt{exec.security}, \texttt{exec.ask}, \texttt{exec.node}.
  \item Update tool plumbing to respect \texttt{exec.host}.
  \item Add \texttt{/exec} slash command and keep \texttt{/elevated} alias.
\end{itemize}


\subsubsection{Phase 2: approvals store + gateway enforcement}


\begin{itemize}
  \item Implement \texttt{exec-approvals.json} reader/writer.
  \item Enforce allowlist + ask modes for \texttt{gateway} host.
  \item Add output caps.
\end{itemize}


\subsubsection{Phase 3: node runner enforcement}


\begin{itemize}
  \item Update node runner to enforce allowlist + ask.
  \item Add Unix socket prompt bridge to macOS app UI.
  \item Wire \texttt{askFallback}.
\end{itemize}


\subsubsection{Phase 4: events}


\begin{itemize}
  \item Add node → gateway Bridge events for exec lifecycle.
  \item Map to \texttt{enqueueSystemEvent} for agent prompts.
\end{itemize}


\subsubsection{Phase 5: UI polish}


\begin{itemize}
  \item Mac app: allowlist editor, per-agent switcher, ask policy UI.
  \item Node binding controls (optional).
\end{itemize}


\subsection{Testing plan}


\begin{itemize}
  \item Unit tests: allowlist matching (glob + case-insensitive).
  \item Unit tests: policy resolution precedence (tool param → agent override → global).
  \item Integration tests: node runner deny/allow/ask flows.
  \item Bridge event tests: node event → system event routing.
\end{itemize}


\subsection{Open risks}


\begin{itemize}
  \item UI unavailability: ensure \texttt{askFallback} is respected.
  \item Long-running commands: rely on timeout + output caps.
  \item Multi-node ambiguity: error unless node binding or explicit node param.
\end{itemize}


\subsection{Related docs}


\begin{itemize}
  \item \href{/tools/exec}{Exec tool}
  \item \href{/tools/exec-approvals}{Exec approvals}
  \item \href{/nodes}{Nodes}
  \item \href{/tools/elevated}{Elevated mode}
\end{itemize}



\clearpage

% === Source file: refactor/outbound-session-mirroring.md ===
\section{Outbound Session Mirroring Refactor (Issue \#1520)}


\subsection{Status}


\begin{itemize}
  \item In progress.
  \item Core + plugin channel routing updated for outbound mirroring.
  \item Gateway send now derives target session when sessionKey is omitted.
\end{itemize}


\subsection{Context}


Outbound sends were mirrored into the \_current\_ agent session (tool session key) rather than the target channel session. Inbound routing uses channel/peer session keys, so outbound responses landed in the wrong session and first-contact targets often lacked session entries.

\subsection{Goals}


\begin{itemize}
  \item Mirror outbound messages into the target channel session key.
  \item Create session entries on outbound when missing.
  \item Keep thread/topic scoping aligned with inbound session keys.
  \item Cover core channels plus bundled extensions.
\end{itemize}


\subsection{Implementation Summary}


\begin{itemize}
  \item New outbound session routing helper:
  \item \texttt{src/infra/outbound/outbound-session.ts}
  \item \texttt{resolveOutboundSessionRoute} builds target sessionKey using \texttt{buildAgentSessionKey} (dmScope + identityLinks).
  \item \texttt{ensureOutboundSessionEntry} writes minimal \texttt{MsgContext} via \texttt{recordSessionMetaFromInbound}.
  \item \texttt{runMessageAction} (send) derives target sessionKey and passes it to \texttt{executeSendAction} for mirroring.
  \item \texttt{message-tool} no longer mirrors directly; it only resolves agentId from the current session key.
  \item Plugin send path mirrors via \texttt{appendAssistantMessageToSessionTranscript} using the derived sessionKey.
  \item Gateway send derives a target session key when none is provided (default agent), and ensures a session entry.
\end{itemize}


\subsection{Thread/Topic Handling}


\begin{itemize}
  \item Slack: replyTo/threadId -> \texttt{resolveThreadSessionKeys} (suffix).
  \item Discord: threadId/replyTo -> \texttt{resolveThreadSessionKeys} with \texttt{useSuffix=false} to match inbound (thread channel id already scopes session).
  \item Telegram: topic IDs map to \texttt{chatId:topic:} via \texttt{buildTelegramGroupPeerId}.
\end{itemize}


\subsection{Extensions Covered}


\begin{itemize}
  \item Matrix, MS Teams, Mattermost, BlueBubbles, Nextcloud Talk, Zalo, Zalo Personal, Nostr, Tlon.
  \item Notes:
  \item Mattermost targets now strip \texttt{@} for DM session key routing.
  \item Zalo Personal uses DM peer kind for 1:1 targets (group only when \texttt{group:} is present).
  \item BlueBubbles group targets strip \texttt{chat\_*} prefixes to match inbound session keys.
  \item Slack auto-thread mirroring matches channel ids case-insensitively.
  \item Gateway send lowercases provided session keys before mirroring.
\end{itemize}


\subsection{Decisions}


\begin{itemize}
  \item \textbf{Gateway send session derivation}: if \texttt{sessionKey} is provided, use it. If omitted, derive a sessionKey from target + default agent and mirror there.
  \item \textbf{Session entry creation}: always use \texttt{recordSessionMetaFromInbound} with \texttt{Provider/From/To/ChatType/AccountId/Originating*} aligned to inbound formats.
  \item \textbf{Target normalization}: outbound routing uses resolved targets (post \texttt{resolveChannelTarget}) when available.
  \item \textbf{Session key casing}: canonicalize session keys to lowercase on write and during migrations.
\end{itemize}


\subsection{Tests Added/Updated}


\begin{itemize}
  \item \texttt{src/infra/outbound/outbound.test.ts}
  \item Slack thread session key.
  \item Telegram topic session key.
  \item dmScope identityLinks with Discord.
  \item \texttt{src/agents/tools/message-tool.test.ts}
  \item Derives agentId from session key (no sessionKey passed through).
  \item \texttt{src/gateway/server-methods/send.test.ts}
  \item Derives session key when omitted and creates session entry.
\end{itemize}


\subsection{Open Items / Follow-ups}


\begin{itemize}
  \item Voice-call plugin uses custom \texttt{voice:} session keys. Outbound mapping is not standardized here; if message-tool should support voice-call sends, add explicit mapping.
  \item Confirm if any external plugin uses non-standard \texttt{From/To} formats beyond the bundled set.
\end{itemize}


\subsection{Files Touched}


\begin{itemize}
  \item \texttt{src/infra/outbound/outbound-session.ts}
  \item \texttt{src/infra/outbound/outbound-send-service.ts}
  \item \texttt{src/infra/outbound/message-action-runner.ts}
  \item \texttt{src/agents/tools/message-tool.ts}
  \item \texttt{src/gateway/server-methods/send.ts}
  \item Tests in:
  \item \texttt{src/infra/outbound/outbound.test.ts}
  \item \texttt{src/agents/tools/message-tool.test.ts}
  \item \texttt{src/gateway/server-methods/send.test.ts}
\end{itemize}



\clearpage

% === Source file: refactor/plugin-sdk.md ===
\section{Plugin SDK + Runtime Refactor Plan}


Goal: every messaging connector is a plugin (bundled or external) using one stable API.
No plugin imports from \texttt{src/**} directly. All dependencies go through the SDK or runtime.

\subsection{Why now}


\begin{itemize}
  \item Current connectors mix patterns: direct core imports, dist-only bridges, and custom helpers.
  \item This makes upgrades brittle and blocks a clean external plugin surface.
\end{itemize}


\subsection{Target architecture (two layers)}


\subsubsection{1) Plugin SDK (compile-time, stable, publishable)}


Scope: types, helpers, and config utilities. No runtime state, no side effects.

Contents (examples):

\begin{itemize}
  \item Types: \texttt{ChannelPlugin}, adapters, \texttt{ChannelMeta}, \texttt{ChannelCapabilities}, \texttt{ChannelDirectoryEntry}.
  \item Config helpers: \texttt{buildChannelConfigSchema}, \texttt{setAccountEnabledInConfigSection}, \texttt{deleteAccountFromConfigSection},
\end{itemize}

\texttt{applyAccountNameToChannelSection}.
\begin{itemize}
  \item Pairing helpers: \texttt{PAIRING\_APPROVED\_MESSAGE}, \texttt{formatPairingApproveHint}.
  \item Onboarding helpers: \texttt{promptChannelAccessConfig}, \texttt{addWildcardAllowFrom}, onboarding types.
  \item Tool param helpers: \texttt{createActionGate}, \texttt{readStringParam}, \texttt{readNumberParam}, \texttt{readReactionParams}, \texttt{jsonResult}.
  \item Docs link helper: \texttt{formatDocsLink}.
\end{itemize}


Delivery:

\begin{itemize}
  \item Publish as \texttt{openclaw/plugin-sdk} (or export from core under \texttt{openclaw/plugin-sdk}).
  \item Semver with explicit stability guarantees.
\end{itemize}


\subsubsection{2) Plugin Runtime (execution surface, injected)}


Scope: everything that touches core runtime behavior.
Accessed via \texttt{OpenClawPluginApi.runtime} so plugins never import \texttt{src/**}.

Proposed surface (minimal but complete):

\begin{lstlisting}[language=typescript]
export type PluginRuntime = {
  channel: {
    text: {
      chunkMarkdownText(text: string, limit: number): string[];
      resolveTextChunkLimit(cfg: OpenClawConfig, channel: string, accountId?: string): number;
      hasControlCommand(text: string, cfg: OpenClawConfig): boolean;
    };
    reply: {
      dispatchReplyWithBufferedBlockDispatcher(params: {
        ctx: unknown;
        cfg: unknown;
        dispatcherOptions: {
          deliver: (payload: {
            text?: string;
            mediaUrls?: string[];
            mediaUrl?: string;
          }) => void | Promise<void>;
          onError?: (err: unknown, info: { kind: string }) => void;
        };
      }): Promise<void>;
      createReplyDispatcherWithTyping?: unknown; // adapter for Teams-style flows
    };
    routing: {
      resolveAgentRoute(params: {
        cfg: unknown;
        channel: string;
        accountId: string;
        peer: { kind: RoutePeerKind; id: string };
      }): { sessionKey: string; accountId: string };
    };
    pairing: {
      buildPairingReply(params: { channel: string; idLine: string; code: string }): string;
      readAllowFromStore(channel: string): Promise<string[]>;
      upsertPairingRequest(params: {
        channel: string;
        id: string;
        meta?: { name?: string };
      }): Promise<{ code: string; created: boolean }>;
    };
    media: {
      fetchRemoteMedia(params: { url: string }): Promise<{ buffer: Buffer; contentType?: string }>;
      saveMediaBuffer(
        buffer: Uint8Array,
        contentType: string | undefined,
        direction: "inbound" | "outbound",
        maxBytes: number,
      ): Promise<{ path: string; contentType?: string }>;
    };
    mentions: {
      buildMentionRegexes(cfg: OpenClawConfig, agentId?: string): RegExp[];
      matchesMentionPatterns(text: string, regexes: RegExp[]): boolean;
    };
    groups: {
      resolveGroupPolicy(
        cfg: OpenClawConfig,
        channel: string,
        accountId: string,
        groupId: string,
      ): {
        allowlistEnabled: boolean;
        allowed: boolean;
        groupConfig?: unknown;
        defaultConfig?: unknown;
      };
      resolveRequireMention(
        cfg: OpenClawConfig,
        channel: string,
        accountId: string,
        groupId: string,
        override?: boolean,
      ): boolean;
    };
    debounce: {
      createInboundDebouncer<T>(opts: {
        debounceMs: number;
        buildKey: (v: T) => string | null;
        shouldDebounce: (v: T) => boolean;
        onFlush: (entries: T[]) => Promise<void>;
        onError?: (err: unknown) => void;
      }): { push: (v: T) => void; flush: () => Promise<void> };
      resolveInboundDebounceMs(cfg: OpenClawConfig, channel: string): number;
    };
    commands: {
      resolveCommandAuthorizedFromAuthorizers(params: {
        useAccessGroups: boolean;
        authorizers: Array<{ configured: boolean; allowed: boolean }>;
      }): boolean;
    };
  };
  logging: {
    shouldLogVerbose(): boolean;
    getChildLogger(name: string): PluginLogger;
  };
  state: {
    resolveStateDir(cfg: OpenClawConfig): string;
  };
};
\end{lstlisting}


Notes:

\begin{itemize}
  \item Runtime is the only way to access core behavior.
  \item SDK is intentionally small and stable.
  \item Each runtime method maps to an existing core implementation (no duplication).
\end{itemize}


\subsection{Migration plan (phased, safe)}


\subsubsection{Phase 0: scaffolding}


\begin{itemize}
  \item Introduce \texttt{openclaw/plugin-sdk}.
  \item Add \texttt{api.runtime} to \texttt{OpenClawPluginApi} with the surface above.
  \item Maintain existing imports during a transition window (deprecation warnings).
\end{itemize}


\subsubsection{Phase 1: bridge cleanup (low risk)}


\begin{itemize}
  \item Replace per-extension \texttt{core-bridge.ts} with \texttt{api.runtime}.
  \item Migrate BlueBubbles, Zalo, Zalo Personal first (already close).
  \item Remove duplicated bridge code.
\end{itemize}


\subsubsection{Phase 2: light direct-import plugins}


\begin{itemize}
  \item Migrate Matrix to SDK + runtime.
  \item Validate onboarding, directory, group mention logic.
\end{itemize}


\subsubsection{Phase 3: heavy direct-import plugins}


\begin{itemize}
  \item Migrate MS Teams (largest set of runtime helpers).
  \item Ensure reply/typing semantics match current behavior.
\end{itemize}


\subsubsection{Phase 4: iMessage pluginization}


\begin{itemize}
  \item Move iMessage into \texttt{extensions/imessage}.
  \item Replace direct core calls with \texttt{api.runtime}.
  \item Keep config keys, CLI behavior, and docs intact.
\end{itemize}


\subsubsection{Phase 5: enforcement}


\begin{itemize}
  \item Add lint rule / CI check: no \texttt{extensions/\textbf{} imports from \texttt{src/}}.
  \item Add plugin SDK/version compatibility checks (runtime + SDK semver).
\end{itemize}


\subsection{Compatibility and versioning}


\begin{itemize}
  \item SDK: semver, published, documented changes.
  \item Runtime: versioned per core release. Add \texttt{api.runtime.version}.
  \item Plugins declare a required runtime range (e.g., \texttt{openclawRuntime: ">=2026.2.0"}).
\end{itemize}


\subsection{Testing strategy}


\begin{itemize}
  \item Adapter-level unit tests (runtime functions exercised with real core implementation).
  \item Golden tests per plugin: ensure no behavior drift (routing, pairing, allowlist, mention gating).
  \item A single end-to-end plugin sample used in CI (install + run + smoke).
\end{itemize}


\subsection{Open questions}


\begin{itemize}
  \item Where to host SDK types: separate package or core export?
  \item Runtime type distribution: in SDK (types only) or in core?
  \item How to expose docs links for bundled vs external plugins?
  \item Do we allow limited direct core imports for in-repo plugins during transition?
\end{itemize}


\subsection{Success criteria}


\begin{itemize}
  \item All channel connectors are plugins using SDK + runtime.
  \item No \texttt{extensions/\textbf{} imports from \texttt{src/}}.
  \item New connector templates depend only on SDK + runtime.
  \item External plugins can be developed and updated without core source access.
\end{itemize}


Related docs: \href{/tools/plugin}{Plugins}, \href{/channels/index}{Channels}, \href{/gateway/configuration}{Configuration}.


\clearpage

% === Source file: refactor/strict-config.md ===
\section{Strict config validation (doctor-only migrations)}


\subsection{Goals}


\begin{itemize}
  \item \textbf{Reject unknown config keys everywhere} (root + nested), except root \texttt{\$schema} metadata.
  \item \textbf{Reject plugin config without a schema}; don’t load that plugin.
  \item \textbf{Remove legacy auto-migration on load}; migrations run via doctor only.
  \item \textbf{Auto-run doctor (dry-run) on startup}; if invalid, block non-diagnostic commands.
\end{itemize}


\subsection{Non-goals}


\begin{itemize}
  \item Backward compatibility on load (legacy keys do not auto-migrate).
  \item Silent drops of unrecognized keys.
\end{itemize}


\subsection{Strict validation rules}


\begin{itemize}
  \item Config must match the schema exactly at every level.
  \item Unknown keys are validation errors (no passthrough at root or nested), except root \texttt{\$schema} when it is a string.
  \item \texttt{plugins.entries..config} must be validated by the plugin’s schema.
  \item If a plugin lacks a schema, \textbf{reject plugin load} and surface a clear error.
  \item Unknown \texttt{channels.} keys are errors unless a plugin manifest declares the channel id.
  \item Plugin manifests (\texttt{openclaw.plugin.json}) are required for all plugins.
\end{itemize}


\subsection{Plugin schema enforcement}


\begin{itemize}
  \item Each plugin provides a strict JSON Schema for its config (inline in the manifest).
  \item Plugin load flow:
\end{itemize}

\begin{enumerate}
  \item Resolve plugin manifest + schema (\texttt{openclaw.plugin.json}).
  \item Validate config against the schema.
  \item If missing schema or invalid config: block plugin load, record error.
\end{enumerate}

\begin{itemize}
  \item Error message includes:
  \item Plugin id
  \item Reason (missing schema / invalid config)
  \item Path(s) that failed validation
  \item Disabled plugins keep their config, but Doctor + logs surface a warning.
\end{itemize}


\subsection{Doctor flow}


\begin{itemize}
  \item Doctor runs \textbf{every time} config is loaded (dry-run by default).
  \item If config invalid:
  \item Print a summary + actionable errors.
  \item Instruct: \texttt{openclaw doctor --fix}.
  \item \texttt{openclaw doctor --fix}:
  \item Applies migrations.
  \item Removes unknown keys.
  \item Writes updated config.
\end{itemize}


\subsection{Command gating (when config is invalid)}


Allowed (diagnostic-only):

\begin{itemize}
  \item \texttt{openclaw doctor}
  \item \texttt{openclaw logs}
  \item \texttt{openclaw health}
  \item \texttt{openclaw help}
  \item \texttt{openclaw status}
  \item \texttt{openclaw gateway status}
\end{itemize}


Everything else must hard-fail with: “Config invalid. Run \texttt{openclaw doctor --fix}.”

\subsection{Error UX format}


\begin{itemize}
  \item Single summary header.
  \item Grouped sections:
  \item Unknown keys (full paths)
  \item Legacy keys / migrations needed
  \item Plugin load failures (plugin id + reason + path)
\end{itemize}


\subsection{Implementation touchpoints}


\begin{itemize}
  \item \texttt{src/config/zod-schema.ts}: remove root passthrough; strict objects everywhere.
  \item \texttt{src/config/zod-schema.providers.ts}: ensure strict channel schemas.
  \item \texttt{src/config/validation.ts}: fail on unknown keys; do not apply legacy migrations.
  \item \texttt{src/config/io.ts}: remove legacy auto-migrations; always run doctor dry-run.
  \item \texttt{src/config/legacy*.ts}: move usage to doctor only.
  \item \texttt{src/plugins/*}: add schema registry + gating.
  \item CLI command gating in \texttt{src/cli}.
\end{itemize}


\subsection{Tests}


\begin{itemize}
  \item Unknown key rejection (root + nested).
  \item Plugin missing schema → plugin load blocked with clear error.
  \item Invalid config → gateway startup blocked except diagnostic commands.
  \item Doctor dry-run auto; \texttt{doctor --fix} writes corrected config.
\end{itemize}



\clearpage
